% Generated draft: requires mathematical and source review.
\documentclass[11pt]{article}
\usepackage[margin=25mm]{geometry}
\usepackage{fontspec}
\setmainfont{Latin Modern Roman}
\usepackage{amsmath,amssymb,mathtools,unicode-math}
\setmathfont{Latin Modern Math}
\usepackage{enumitem,needspace,xurl,hyperref}
\setlength{\parindent}{0pt}
\setlength{\parskip}{5pt}
\setlength{\emergencystretch}{3em}
\newcommand{\N}{\mathbb{N}}
\newcommand{\Z}{\mathbb{Z}}
\newcommand{\Q}{\mathbb{Q}}
\newcommand{\R}{\mathbb{R}}
\newcommand{\C}{\mathbb{C}}
\newcommand{\F}{\mathbb{F}}
\newcommand{\A}{\mathbb{A}}
\newcommand{\PP}{\mathbb{P}}
\newcommand{\E}{\mathbb{E}}
\newcommand{\eps}{\varepsilon}
\newcommand{\dd}{\,\mathrm{d}}
\begin{document}
\section*{UnsolvedMath ID: 11000093}
\textbf{Source problem code:} catalogue entry
\subsection{Definitions and mathematical statement}
Conjecture 4.8 on PDF page 79 states that every component of the source's hyperelliptic locus \(\mathcal H_g^c\) is simply connected. The spaces \(\mathcal H_g^c\), \(\mathcal H_g\), their marking conventions and the superscript \(c\) are not defined in this extract. The source says that in genus \(2\) there is one component, equal to \(\mathfrak h_2\), so the assertion holds there. This is a different locally numbered Conjecture 4.8 from any identically numbered item in another chapter.

The source supplies a group-theoretic equivalent. Fix a hyperelliptic involution \(\sigma\) of the closed genus-\(g\) surface, with \(\Delta_g\) the hyperelliptic mapping class group, conventionally its centralizer, and \(T_g\) the Torelli group. Set
\[
 T\Delta_g=\Delta_g\cap T_g
       =\ker\bigl(\Delta_g\to\mathrm{Sp}(H)\bigr),
 \qquad H=H_1(\Sigma_g;\Z).
\]
The source writes the symplectic target \(\mathrm{Sp}_g(\Z)\); its indexing is expressed here without changing the underlying homology action. It places \(T\Delta_g\) inside the Johnson subgroup
\[
 K_g=\ker\bigl(\tau:T_g\to\bigwedge^3H/H\bigr).
\]
The Johnson map and embedding of \(H\) in the displayed quotient need their source conventions.

Let \(N_g\) be the normal subgroup of \(T\Delta_g\) generated by the Dehn twists about separating simple closed curves invariant under \(\sigma\). The source identifies, for a component \(\mathcal H_{g,\alpha}\),
\[
 \pi_1(\mathcal H_{g,\alpha},*)\cong T\Delta_g,
 \qquad
 \pi_1(\mathcal H_{g,\alpha}^c,*)\cong T\Delta_g/N_g.
\]
It attributes the second quotient to van Kampen, identifying twist conjugacy classes with loops around components of the divisor \(\mathcal H_g^{c,\rm red}\). Accordingly its equivalent conjectural statement is \(T\Delta_g=N_g\), expressed in the source as generation by those twist conjugacy classes. The ambient group for normal closure and the geometric identifications must be checked against the missing locus definitions.

For genus \(3\), the source says the conjecture directly implies the known generation of \(T_3\) by \(35=36-1\) bounding-pair elements. It calls \(35\) the rank of \(H_1(\mathcal J_3,\mathcal H_3)\), and says oriented inverse images of generators of \(\pi_1(\mathcal J_3,\mathcal H_3)\) generate \(T_3\). These comma expressions are retained literally: relative homotopy notation, a complement or a lost punctuation mark would give different objects, and the lifting map is absent. Footnote 10 is not supplied. Finally, the source motivates studying hyperelliptic loci because they form the period-map branch locus; it supplies no complete period-map framework in this extract.
\subsection{Short English statement}
The conjecture says each specified hyperelliptic-locus component is simply connected. The source reformulates it as generation of the hyperelliptic Torelli group by conjugacy classes of symmetric separating twists, and gives genus-two and genus-three context. The precise loci and some homology notation are missing.
\subsection{Sources}
\begin{itemize}
\item Local UnsolvedMath registry, permanent ID \(11000093\), exact source code \texttt{catalogue entry}, Conjecture 4.8, PDF page 79. No individual proposer or proposed year is supplied.
\item Farb, edited volume \emph{Problems on Mapping Class Groups and Related Topics}: \url{https://www.math.uchicago.edu/~farb/papers/mcgbook.pdf}. The volume has 329 explicitly typeset numbered items across chapters, with repeated local numbering disambiguated by PDF page. The extract contains a footnote 10 marker without text.
\item Extraction pdf-font-structure, accessed 2026-07-29. Source evidence NEEDS\_REVIEW presents an open problem and requires a current primary-literature audit. Registry status open and classification OPEN-TRIAGE; the research summary explicitly cannot verify status. No separate full references for the group-generation facts or period-map identifications are supplied.
\item Dataset attribution: Ulam AI / UnsolvedMath Contributors, CC BY 4.0. Source rights NEEDS\_REVIEW: publicly accessible author/publisher copy, redistribution terms requiring review.
\end{itemize}
\end{document}
